\documentclass[10pt]{article}
\usepackage[T1]{fontenc}
\usepackage{lmodern,amsmath,amssymb,amsthm,microtype}
\usepackage[margin=0.8in]{geometry}
\usepackage[colorlinks=true,urlcolor=blue]{hyperref}
\newtheorem{theorem}{Theorem}
\newtheorem{lemma}{Lemma}
\newtheorem{corollary}{Corollary}
\title{Squares Whose Equal-Length Halves Differ by a Square}
\author{Qichao Wang\\{\small Hebei University of Technology}\\[0.3em]
Daoyu Dong\\{\small University of Electronic Science and Technology of China}}
\date{October 26, 2026}
\hypersetup{pdftitle={All-Index Existence and Exact Minima for OEIS A236383},pdfauthor={Qichao Wang and Daoyu Dong}}
\begin{document}
\maketitle
\begin{abstract}
For every positive integer $n$, we construct a square whose decimal
expansion is the concatenation of equal-length positive integers
$x,y$ with $y-x=n^2$. This proves the existence conjecture in OEIS
A236383. An odd decimal length in one prescribed residue class
modulo $22$ suffices. Its size is bounded by the number of digits
of $4(n+1)^2$ plus $21$. We then give an exact finite minimum
formula using square roots of $n^2$ modulo $10^d+1$, with explicit
prime-power root lists and the Chinese remainder theorem.
\end{abstract}

\section{A uniform elementary construction}
OEIS A236383 \cite{oeis} defines $a(n)$ as the least positive $k$
such that $k^2$ concatenates two positive equal-length integers
$x,y$ with $y=x+n^2$. Thus, for some $d\ge1$,
\begin{equation}\label{eq:def}
 k^2=10^d x+y=(10^d+1)x+n^2,\qquad
 10^{d-1}\le x\le y<10^d.
\end{equation}
The entry conjectures existence for every $n\ge1$.

\begin{lemma}\label{lem:eleven}
For every positive odd $d$, the integer
$C_d=(10^d+1)/11$ satisfies $C_d\equiv d\pmod{11}$.
\end{lemma}
\begin{proof}
Since $d$ is odd, the binomial expansion of $(-1+11)^d+1$
is congruent to $11d$ modulo $121$. Divide by $11$.
\end{proof}

Let $D$ be the number of decimal digits of $4(n+1)^2$, and let
$d_*$ be the least integer $d\ge D$ satisfying
\begin{equation}\label{eq:length}
 d\text{ odd},\qquad 5d+2n\equiv0\pmod{11}.
\end{equation}
The two conditions specify exactly one residue class modulo $22$,
so $D\le d_*\le D+21$.

\begin{theorem}\label{thm:witness}
Put $d=d_*$, $B=10^d$, $M=B+1$, and $C=M/11$. Then
\begin{equation}\label{eq:construction}
 k=n+5C,\qquad x=\frac{5(2n+5C)}{11},\qquad y=x+n^2
\end{equation}
are positive integers satisfying \eqref{eq:def}. In particular
$a(n)$ exists at every positive index, and $a(n)\le n+5C$.
\end{theorem}
\begin{proof}
Lemma~\ref{lem:eleven} and \eqref{eq:length} show that
$2n+5C\equiv0\pmod{11}$, proving integrality of $x$.
Moreover,
\[
 k^2-n^2=5C(2n+5C)=Mx.
\]
It remains to check both length inequalities, which congruence
alone would not guarantee. The expression for $x$ is
\[
 x=\frac{25}{121}M+\frac{10}{11}n.
\]
Since $25/121>1/10$, it gives $x>B/10=10^{d-1}$.
Also $B\ge10^D>4(n+1)^2$, and therefore
\[
 \frac{96}{121}M>(n+1)^2>
 n^2+\frac{10}{11}n+1.
\]
Rearranging gives $y=x+n^2<M-1=B$. Thus $x,y$ both have
exactly $d$ digits and all conditions in \eqref{eq:def} hold.
There is no exception when $11\mid n$: the prescribed residue
class for $d$ then simply has $11\mid d$.
\end{proof}

\section{The exact least value at every index}
For $1\le d\le d_*$ put $B_d=10^d$, $M_d=B_d+1$, and let
\[
 R_d(n)=\{0\le r<M_d:r^2\equiv n^2\pmod{M_d}\}.
\]
If $B_d-1-n^2<10^{d-1}$, set $\mathcal K_d(n)=\varnothing$.
Otherwise define the exact integer endpoints
\begin{align*}
 L_d(n)&=\left\lceil\sqrt{M_d10^{d-1}+n^2}\right\rceil,\\
 H_d(n)&=\left\lfloor\sqrt{M_d(B_d-1-n^2)+n^2}\right\rfloor,\\
 \mathcal K_d(n)&=R_d(n)\cap[L_d(n),H_d(n)].
\end{align*}
\begin{theorem}\label{thm:min}
For every positive $n$,
\[
 \boxed{a(n)=\min\bigcup_{d=1}^{d_*}\mathcal K_d(n).}
\]
The union is nonempty and all root lists are given by the finite
prime-power expressions below. The evaluator can stop at the first
$d$ with $\mathcal K_d(n)\ne\varnothing$.
\end{theorem}
\begin{proof}
For fixed $d$, \eqref{eq:def} is equivalent to the displayed
congruence and the two bounds on $k^2$. Any valid $k$ satisfies
\[
 k^2\le M_d(B_d-1-n^2)+n^2<B_d^2,
\]
so $k<M_d$; no translates of the root residues are missing.
Theorem~\ref{thm:witness} supplies a candidate at $d=d_*$.
For any admissible length $d$, \eqref{eq:def} implies that $k^2$
has $2d$ digits, so $k$ has exactly $d$ digits. A larger $d$ can
therefore never give a smaller root $k$. This also proves the
stopping rule and excludes all lengths beyond $d_*$.
\end{proof}

\section{Explicit local root lists}
Every $M_d$ is odd. For an odd prime power $p^e\parallel M_d$,
write $n=p^v u$ with $p\nmid u$. If $2v\ge e$, the complete
list of roots modulo $p^e$ is
\begin{equation}\label{eq:zero}
 r=j p^{\lceil e/2\rceil},\qquad 0\le j<p^{\lfloor e/2\rfloor}.
\end{equation}
If $2v<e$, the complete list is
\begin{equation}\label{eq:nonzero}
 r=p^v\bigl(\varepsilon u+j p^{e-2v}\bigr)\pmod{p^e},
 \qquad\varepsilon\in\{-1,1\},\quad0\le j<p^v.
\end{equation}
To verify completeness, in the first case $r^2\equiv0$ is
equivalent to $p^{\lceil e/2\rceil}\mid r$. In the second case
any root must have valuation $v$; after division by $p^{2v}$,
the equation is $z^2\equiv u^2\pmod{p^{e-2v}}$.
The factors $z-u,z+u$ cannot both be divisible by the odd prime
$p$, since $p\nmid2u$. Thus $z\equiv u$ or $-u$ modulo that
prime power; its lifts give precisely \eqref{eq:nonzero}.
The Chinese remainder theorem combines these local lists to form
all of $R_d(n)$. This avoids a scan of consecutive $k$ values.

The formula is an exact finite arithmetic expression involving
factorization, explicit root lists, integer square roots and a
finite minimum. We do not assert an elementary closed expression
without that final minimum. Existence and the length bound are
proved for every $n$, rather than inferred from the finite tests.

\section{Independent verification}
The exact root evaluator reproduces all 245 official b-file terms.
An independent exhaustive scan of $1\le k\le1000000$ finds 659
attained indices and their minima; all agree with the evaluator,
including all 245 published indices. Further checks compare the
prime-power and composite root lists with direct residue enumeration.
The explicit construction is checked through $n=1000$, at powers
$11^1,\ldots,11^{20}$, and at a 101-digit index. The accompanying
300-term table extends the official published range. These checks
validate the implementations of the all-index results above.

\begin{thebibliography}{9}
\bibitem{oeis} The OEIS Foundation Inc.,
\href{https://oeis.org/A236383}{A236383} and its
\href{https://oeis.org/A236383/b236383.txt}{245-term b-file},
accessed October 26, 2026.
\end{thebibliography}
\end{document}
